\section{The development loop}
\label{sec:ch02-dev-loop}

A schema you did not write by hand is a schema you have to go and look at. That
is the one genuine cost of the implementation-first style, and it shapes how you
work day to day: the loop is edit, restart, look at what the schema became.

\index{HotChocolate!eager initialisation}%
The restart is better than it used to be. HotChocolate~16 builds the schema
eagerly during application startup rather than lazily on the first
request~\autocite{chillicream2026migrate16}. If you have written HotChocolate
before, this is the change most likely to alter your habits, and it is worth
being concrete about the difference. Under lazy initialisation a schema mistake
was a broken response to whoever sent the first query, which in practice meant
you found it in the IDE and sometimes in staging. Now the process fails to start.
A container with a bad schema never becomes healthy, which turns a class of
error that used to reach a client into a deployment that visibly did not happen.
If you had \code{InitializeOnStartup()} in an older service, delete it; the
behaviour is the default now.

\index{schema!export}%
The second habit is exporting the schema to a file. The template wires this up
already, through the last line of \code{Program.cs}:

\begin{minted}{csharp}
app.RunWithGraphQLCommands(args);
\end{minted}

That call gives the executable a command line as well as a web server, so the
same project that serves the graph can print it:

\begin{minted}{text}
dotnet run --project src/Mosaic.Api -- \
    schema export --output schema/mosaic.graphql
\end{minted}

The exporter also drops a \code{mosaic-settings.json} next to the SDL, which you
almost certainly do not want in the repository; Mosaic gitignores it.

Mosaic commits the schema itself. That is the part I would push hardest on if I
were reviewing somebody else's HotChocolate service, because it converts an
invisible change into a visible one. Rename a resolver method and the schema changes;
nothing in the C\# diff says so, and a reviewer reading the pull request has to
know HotChocolate's naming rules well enough to notice. With the schema in the
repository, the same pull request carries a diff that says \code{price} became
\code{unitPrice} in as many words. Mosaic's verify script regenerates the file
and fails if it differs from the committed copy, so the two cannot drift apart
quietly.

That check earns its place in a way I did not anticipate. Exported SDL turns out
to include things nobody typed:

\begin{graphqlsdl}
type Query {
  products: [Product!]! @cost(weight: "10")
  productById(id: ID!): Product @cost(weight: "10")
}
\end{graphqlsdl}

\index{HotChocolate!cost analysis}%
The \code{@cost} directives come from HotChocolate~16's cost analyzer, which is
one of the defaults \code{AddGraphQLServer()} applies, and the weights are
assigned by the server rather than by you. Which fields get one is not arbitrary,
and working out the rule took me an afternoon that I will save you: an
asynchronous resolver gets a default weight, and so does a resolver supplied as
a delegate, because a delegate is opaque to the analyzer. A resolver written as
a method on a class, which HotChocolate compiles and can inspect, gets no weight
at all when it is synchronous and takes nothing but field arguments. I found
this because two of the sample projects from
section~\ref{sec:ch02-two-ways} refused to produce identical schemas, and the
whole difference was one resolver written as a lambda.

\index{Nitro!schema browsing}%
For looking at the schema rather than diffing it, Nitro is already running at
the endpoint. Create a document, point it at
\code{http://localhost:5100/graphql}, and the Schema tab renders the types while
the Schema Definition tab shows raw SDL. It is the fastest way to answer
\enquote{did that field come out with the name I expected}, which in a
generated schema is a question you ask a lot.

So the loop is: change a type, restart and let startup tell you if the schema is
broken, export and read the diff if you meant to change the contract, and open
Nitro when you want to see the shape rather than the text.

\begin{figure}[htbp]
  \centering
  % The edit and check cycle, with the SDL snapshot hanging off it. Referenced
% from chapter 02. Bare tikzpicture: the figure environment, caption and label
% live at the call site. Uses only arrows.meta and positioning.
\begin{tikzpicture}[
    font=\small,
    stage/.style={draw, rounded corners=2pt, minimum width=27mm,
                  minimum height=12mm, align=center},
    cmd/.style={draw, rounded corners=2pt, minimum width=38mm,
                minimum height=12mm, align=center},
    art/.style={draw, thick, rounded corners=2pt, minimum width=34mm,
                minimum height=12mm, align=center},
    flow/.style={-{Stealth[length=2mm]}, thick},
    note/.style={font=\scriptsize\itshape, align=center}
  ]

  \node[stage] (edit)    at (-2.9, 1.5) {edit a\\C\# type};
  \node[stage] (restart) at ( 2.9, 1.5) {restart\\the service};
  \node[stage] (schema)  at ( 2.9,-1.5) {the schema\\changes};
  \node[stage] (probe)   at (-2.9,-1.5) {re-introspect\\in Postman};

  \draw[flow] (edit)    -- (restart);
  \draw[flow] (restart) -- (schema);
  \draw[flow] (schema)  -- (probe);
  \draw[flow] (probe)   -- (edit);

  \node[note, anchor=south] at (0,2.25)
    {the schema is built eagerly at startup:\\
     a mistake fails here, not on the first query};

  \node[note, anchor=north] at (-2.9,-2.25)
    {introspection answers\\only in Development};

  % The branch that turns a schema change into a diff.
  \node[cmd] (export)   at (2.9,-3.5) {\texttt{dotnet run}\\\texttt{-- schema export}};
  \node[art] (snapshot) at (2.9,-5.4) {committed SDL\\snapshot};

  \draw[flow] (schema) -- (export);
  \draw[flow] (export) -- (snapshot);

  \node[note, anchor=north] at (2.9,-6.1)
    {this is what makes a schema\\change show up in a diff};

\end{tikzpicture}

  \caption{The cycle this section describes. The two branches off it are what
  make a generated schema manageable: startup fails loudly when the schema is
  wrong, and the committed SDL file turns a change in the contract into
  something a reviewer can read.}
  \label{fig:ch02-dev-loop}
\end{figure}

There is one more tool in the loop, and it is the one the rest of this book
leans on hardest.
